
\subsubsection{Корректировка параметров волн за счет трансформации на мелководье}
\label{sec:seashoaring}

В прошлом разделе мы получили параметры описывающие состояние волны на открытой, глубокой воде.
Дальше эта \enquote{глубоководная} волна распространяется к берегу по лучу, повторяющему направление ветра.
В какой-то момент она начнет взаимодействовать с неоднородным рельефом дна, что приведет к перераспределению энергии, росту высоты волны на мелководье (эффекту шоалинга) и, в конечном счете, к ее разрушению в момент, когда высота станет слишком большой относительно глубины под собой. 

Необходимость учета шоалинга определяется тем, что он радикально меняет высоту волн на подходе к берегу, а, следовательно, и переносимую волной энергию.
В определенных условиях, волна, которая в глубокой воде кажется умеренной, на мелководье может вырасти в 1.5-2 раза и именно в таком виде воздействовать на прибрежную территорию.
В результате, если использовать только \enquote{глубоководный} $H_s$ из SMB (формула \eqref{eq:smb_hs}) без учета трансформации на мелководье, мы рискуем систематически недооценить прибрежные волновые нагрузки, а значит -- энергию, определяющую интенсивность эрозии, переотложения наносов и разрушения инфраструктуры, что прямо приведет к занижению оценок прибрежной уязвимости~\cite{Zhang2026EffectDirectionalityExtreme}.

Одним из вариантов математического описания изменения высоты волны при ее распространении в мелкой воде является закон Грина, в основе которого лежит идея сохранения потока энергии~\cite{jay_1991_greens_law}.
Гидротехнический опыт доказывает, что закон Грина дает хорошие результаты для случаев когда волна не разрушается за счет потери энергии (вызванными вязкостью, трением, турбулентностью и~т.\,п.)~\cite{greens_law_solitary_waves}.

С учетом заданных ограничений поток энергии вдоль луча можно считать постоянным, что при выполнении следующих условий:
\begin{itemize}
    \item волна длинная (длина волны намного больше глубины ($\lambda \gg h$));
    \item глубина $h(x)$ меняется медленно вдоль луча;
    \item фронт волны примерно параллелен изобатам;
    \item ширина луча постоянна.
\end{itemize}

Это позволяет принять следующее соотношение между высотой волны $H$ и глубиной $h$:
\begin{equation}\label{eq:green_low}
    H\,h^{1/4} = \text{const}.
\end{equation}

Из выражения~\eqref{eq:green_low} высота волны $H$ зависит от глубины $h$ как:
\begin{equation}
    H(h) \propto h^{-1/4}.
\end{equation}

Одновременно длина волны $\lambda$ при этом изменяется так, что выполняется:
\begin{equation}
    \frac{\lambda}{\sqrt{g\,h}} = \text{const},
\end{equation}
то есть в мелкой воде $\lambda \propto \sqrt{h}$, а период $T$ остается постоянным.

При переходе от глубокой воды к мелководью высота волны в модели корректируется с помощью коэффициента шоалинга $K_s$ -- безразмерного множителя, определяющего кратность изменения высоты волны за счет изменения глубины, при прочих равных условиях:
\begin{equation}
    K_s(h) = \frac{H(h)}{H_0},
\end{equation}
где $H_0$ -- высота волны в глубокой воде, 
$H(h)$ -- высота при глубине $h$.

Так при $K_s > 1$ волна будет усиливаться при входе на мелководье (чистый шоалинг), а при $K_s \approx 1$ -- глубина еще достаточно велика и дно почти не влияет на высоту волны.

Для длинных волн в мелкой воде при плавной батиметрии и постоянной ширине луча закон Грина дает:
\begin{equation}
    H(h)\,h^{1/4} = \text{const} \quad\Rightarrow\quad
    \frac{H(h)}{H_0} = \left(\frac{h_0}{h}\right)^{1/4}.
\end{equation}

Отсюда коэффициент шоалинга между глубокой водой ($h_{\text{deep}}$) и прибрежной точкой расчета ($h_{\text{point}}$) может быть найден как:
\begin{equation}
    K_s = \left(\frac{h_{\text{deep}}}{h_{\text{point}}}\right)^{1/4},
    \qquad
    H_{\text{shoaled}} = H_{\text{offshore}} \cdot K_s.
\end{equation}

Так что для любой глубины $h$ вдоль луча:
\begin{equation}\label{eq:green_low_H}
    H(h)=H_{offshore}\left(\frac{h_{deep}}{h}\right)^{1/4}.
\end{equation}

Из формулы~\eqref{eq:green_low_H} очевидно, что при нулевой глубине $h = 0$ она либо приводит к дегенерации волны, так как она выходит за пределы применимости закона Грина (волна на берегу разрушается).
В формуле~\eqref{eq:green_low_H} волна еще рассматривается в воде, а не как уже \enquote{сошедший} на берег прибой.

Таким образом, численное значение параметра $h$ следует взять из батиметрии для максимально близкой к берегу точки с положительной глубиной по направлению рассматриваемого луча.
Тем самым формула~\eqref{eq:green_low_H} описывает поведение на \enquote{подходах} к прибойной зоне, а не на самой линии берега.

Тем не менее, прямое вычисление высоты волны по формуле~\eqref{eq:green_low_H} сохраняет риск получения завышенных значений, так как оно не учитывает изменение рельефа морского дна и возможность ее разрушения до подхода к берегу.

Решить эту проблему можно, добавив в расчет критерий разрушения волны, определяемый соотношением:
\begin{equation}\label{eq:break_wave_eq}
    H(h) \ge \gamma_b\,h,
\end{equation}
где $\gamma_b$ -- безразмерный коэффициент разрушения.

Фактически выражение~\eqref{eq:break_wave_eq} вводит некоторую предельную крутизну волны, ограничивая максимально возможный рост волны и сохраняя общую устойчивость решения.
В литературе $\gamma_b = H_b/h_b$ называют "breaker index", определяемый по данным натурных или лабораторных наблюдений.
Коэффициент $\gamma_b$ лежит примерно в диапазоне $0.55-0.78$, в зависимости от крутизны дна и параметров волны.
Значение 0.55 соответствует нижней границе диапазона и подходит для пологих берегов с нерегулярным волнением, тогда как значение 0.78 традиционно используется в Shore Protection Manual~\cite{cerc_spm_1984} для монохроматических волн. 

Для учета трансформации ветровых волн при подходе к берегу модель  
использует одномерный батиметрический профиль вдоль луча, совпадающий  
с направлением распространения волны. 
Логика контроля критической высоты в предлагаемом методе состоит в последовательном поиске по профилю глубин первой точки, где рост волны по шоалингу становится \enquote{слишком большим} относительно глубины, и дальше высота волны жестко ограничивается этим пределом.

Сначала для каждой глубины $h$ вдоль луча (от моря к берегу) считается высота, которая волна \enquote{хотела бы} иметь при чистом шоалинге по закону Грина:
\begin{equation}\label{eq:green_low_result}
    H(h) = H_{\text{offshore}} \left(\frac{h}{h_{\text{deep}}}\right)^{1/4},
\end{equation}
где в качестве $H_{\text{offshore}}$ берется значимая высота в глубокой воде из SMB, 
$h_{\text{deep}}$ -- максимальная глубина на профиле. 

В результате мы получаем идеализированную \enquote{кривую роста} высоты волны от глубины дна.
После чего для каждой высоты проверяется условие разрушения~\eqref{eq:break_wave_eq}.
Первая глубина $h_i$, для которой условие выполняется, фиксируется как глубина разрушения $h_{\text{break}}$.

Предельная высота разрушившейся волны задается условием:
\begin{equation}
    H_{\text{break}} = \gamma_b\,h_{\text{break}}.\label{eq:break_cond}
\end{equation}

После этого итоговая прибрежная высота волны определяется как минимальная из них:
\begin{equation}
    H_{\text{near}} = \min\!\left( H_{\text{shoaled}},\, H_{\text{break}} \right).
\end{equation}

В результате, если чистый шоалинг~\eqref{eq:green_low_result} дает высоту $H_{\text{shoaled}} \leq H_{\text{break}}$~\eqref{eq:break_cond}, значит волна не успевает «дорасти» до критического состояния и прибрежная высота определяется только законом Грина.
В противном случае, где-то по пути волна должна разрушиться, а значит ее рост ограничивается $H_{\text{break}}$.

Следует отметить, что в реальной прибрежной зоне действуют и процессы, которые потенциально могут препятствовать шоалингу уменьшить высоту волн.
Так, после точки разрушения волны часть энергии переходит в турбулентность и высота дальше к берегу может уменьшаться (бор, прибойный вал).
Морская растительность, пористые рифы и искусственные сооружения могут частично гасить волну, за счет повышенного донного трения, так что прибрежная высота оказывается меньше \enquote{чистой} шоалинговой~\cite{Magdalena2025waveShoalingPorous}.

Следует отдельно отметить, что описанная схема трансформации основывается на ряде упрощающих допущений.  
Во-первых, батиметрия рассматривается вдоль одномерного луча, без учета дифракции и кривизны изобат, что приемлемо для гладких и относительно прямолинейных участков берега. 
Во-вторых, коэффициент шоалинга и критерий разрушения задаются в виде фиксированных степенных и линейных зависимостей, не учитывающих спектральную структуру волнения и направленное распределение энергии. 
Тем не менее, общее решение остается допустимым, потому что выбранные упрощения согласованы с задачей, масштабом и доступными данными, так как модель преследует цель получить прибрежную волну разумного масштаба из ветра и батиметрии, а не точное фазоразрешенное поле волн.
